\documentclass[11pt]{article}
\usepackage[a4paper,margin=25mm]{geometry}
\usepackage{amsmath,amssymb,amsthm,mathtools,hyperref}
\newtheorem{theorem}{Theorem}\newtheorem{lemma}{Lemma}
\newcommand{\R}{\mathbb R}\newcommand{\C}{\mathbb C}
\newcommand{\PF}{\operatorname{PF}}\newcommand{\dd}{\,d}
\title{Completion of the original Montgomery--Weil finite-part identity}
\author{Split-Zero research continuation}\date{19 September 2026}
\begin{document}\maketitle
\noindent The original Six Readers files are preserved unchanged. This note
completes the scalar identity left open after
\texttt{eq:poisson-pole-complete} in
\texttt{proofs/B\_guinand\_weil\_montgomery\_interfaces.tex}, and in the
Montgomery paragraph of
\texttt{satellites/10\_literature\_foundations.tex}.
It independently derives the new incoming note
\texttt{03\_MONTGOMERY\_ARCHIMEDEAN\_IDENTITY.md} from the literal test and
finite-part convention in those complete sources. The distinct false local
Guinand--Weil residual equality is corrected in the companion derivation;
it is not used anywhere in this proof.

\section{The original parameters, transform and domain}
Throughout retain
\begin{equation}\tag{MA1}\begin{gathered}
 1<\sigma_0<2,\quad t\in\R,\quad x\ge1,\quad U=\log x,\quad
 a=\sigma_0-\tfrac12,\\
 F(v)=e^{-a|v-U|}e^{-it(v-U)},\quad
 C=F(0)=x^{1/2-\sigma_0+it},\\
 z=1-\sigma_0+it,\quad w=\sigma_0-it=1-z,\quad
 p=\sigma_0+it,\quad A=(2\pi)^{-1}.
\end{gathered}\end{equation}
For every positive base, a complex power means $y^s=\exp(s\log y)$ with
the real natural logarithm. The digamma function $\psi=\Gamma'/\Gamma$
is meromorphic, not a chosen branch of $\log\Gamma$.
Weil's real-place kernel and the retained finite part \cite{Weil} are
\begin{equation}\tag{MA2}\begin{gathered}
 K(v)=\frac{e^{|v|/2}}{|e^v-e^{-v}|},\\
 \PF(FK)=\lim_{\Lambda\to\infty}
 \left\{\int_\R(1-e^{-\Lambda|v|})F(v)K(v)\dd v-C\log\Lambda\right\}.
\end{gathered}\end{equation}
Indeed $|v|K(v)\to1/2$ and $F(v)\to C$ as $v\to0$, so the original
subtraction $2\beta(0)\log\Lambda$ has $\beta(0)=C/2$. In particular its
coefficient is $C$, not $2C$.

The test $F$ is continuous. On the two intervals separated by $U$ its
derivative is $(a-it)F$ and $(-a-it)F$, respectively. The average derivative
value at the single jump is $-it$, since $F(U)=1$. These are exactly the
first-kind jump data of Weil's test convention. For every
$0<b<\sigma_0-1$,
\[
 |F(v)|+|F'(v)|\le(1+a+|t|)e^{aU}e^{-a|v|}
 =O(e^{-(1/2+b)|v|}).
\]
Thus both literal conditions (A) and (B) of the recorded Weil source hold,
including when $U=0$.
For $1-\sigma_0<\Re s<\sigma_0$, absolute convergence permits $v=U+u$:
\begin{equation}\tag{MA3}\begin{aligned}
 \Phi(s)&:=\int_\R F(v)e^{(s-1/2)v}\dd v\\
 &=x^{s-1/2}\left[
 \int_0^\infty e^{-[a-(s-1/2-it)]u}\dd u
 +\int_0^\infty e^{-[a+(s-1/2-it)]u}\dd u\right]\\
 &=\frac{2a x^{s-1/2}}{a^2-(s-1/2-it)^2}.
\end{aligned}\end{equation}
The strip contains $s=0,1$ and the entire closed critical strip. Factoring
the two original denominators yields
\begin{equation}\tag{MA4}
 \Phi(0)=-\frac{2a x^{-1/2}}{zp},\qquad
 \Phi(1)=\frac{2a x^{1/2}}{(\sigma_0-1+it)(\sigma_0-it)}.
\end{equation}
There is no zero denominator on (MA1): $-1<\Re z<0$, $1<\Re w,\Re p<2$.

\section{The finite part as an ordinary, convergent regularized integral}
On $v>0$ write
\[
 J(v)=[F(v)+F(-v)]K(v).
\]
The exact kernel is $K(v)=e^{-v/2}/(1-e^{-2v})$.
It is $1/(2v)+O(1)$ near zero, while
$F(v)+F(-v)=2C+O(v)$ by local Lipschitz continuity, including at $U=0$.
Consequently $J(v)=C/v+O(1)$. For $v\ge\max(1,U)$, the test has decay
$O_{\sigma_0,t,x}(e^{-av})$ and $K(v)=O(e^{-v/2})$, giving
$J(v)=O_{\sigma_0,t,x}(e^{-\sigma_0v})$.
Therefore the function
\begin{equation}\tag{MA5}
 H(v)=J(v)-C\frac{e^{-v}}v
\end{equation}
is absolutely integrable on $(0,\infty)$.
For finite $\Lambda>0$, elementary differentiation in $\Lambda$ and the
value zero at $\Lambda=0$ prove
\[
 \int_0^\infty\frac{e^{-v}-e^{-(1+\Lambda)v}}v\dd v=\log(1+\Lambda).
\]
It follows exactly that
\begin{equation}\tag{MA6}
 \int_\R(1-e^{-\Lambda|v|})FK\dd v-C\log\Lambda
 =\int_0^\infty(1-e^{-\Lambda v})H(v)\dd v
  +C\log(1+\Lambda^{-1}).
\end{equation}
Dominated convergence with the proved $L^1$ function $H$ now gives
\begin{equation}\tag{MA7}
 \boxed{\PF(FK)=\int_0^\infty
 \left\{[F(v)+F(-v)]K(v)-C\frac{e^{-v}}v\right\}\dd v.}
\end{equation}
This is an exact linear relation between the original finite-part functional
and an ordinary regularized integral. The subtracted function and its
constant have both been retained. The same estimates are uniform when
$(\sigma_0,t,x)$ varies over compact subsets of the domain (MA1), so this
functional is continuous there, including the relative endpoint $x=1$.

\section{Direct integration on the two original intervals}
The positive kernel has the geometric expansion
\begin{equation}\tag{MA8}
 K(v)=\sum_{m=0}^\infty e^{-(2m+1/2)v},\qquad v>0.
\end{equation}
For a fixed $m\ge0$ put $\lambda_m=2m+1/2$. The actual test on the two
intervals is
\begin{equation}\tag{MA9}\begin{aligned}
 F(v)&=Ce^{(a-it)v}&& (0\le v\le U),\\
 F(v)&=Cx^{2a}e^{-(a+it)v}&& (v\ge U),\\
 F(-v)&=Ce^{-(a-it)v}&& (v\ge0).
\end{aligned}\end{equation}
In particular the first interval is empty when $x=1$, and remains finite
when $m=0$ and $\Re z<0$. Direct integration gives
\begin{equation}\tag{MA10}\begin{aligned}
 I_m&:=\int_0^\infty[F(v)+F(-v)]e^{-\lambda_m v}\dd v\\
 &=C\frac{1-e^{-(2m+z)U}}{2m+z}
   +Cx^{2a}\frac{e^{-(2m+p)U}}{2m+p}
   +\frac C{2m+w}\\
 &=C\left(\frac1{2m+w}+\frac1{2m+z}\right)
       -\frac{2a x^{-2m-1/2}}{(2m+z)(2m+p)}.
\end{aligned}\end{equation}
Here $Cx^{-(2m+z)}=Cx^{2a-(2m+p)}=x^{-2m-1/2}$ and $p-z=2a$.
No integral of $e^{-(2m+z)v}$ over the infinite half-line was used.

For finite $\Lambda>2$ all integrations and (MA8) may be exchanged in the
regulated integral. To prove this for complex $F$, apply Tonelli to the
absolute values: the dominating function is
$(1-e^{-\Lambda v})(|F(v)|+|F(-v)|)K(v)$, bounded near zero and
exponentially integrable at infinity by the estimates preceding (MA5).
The same integration as (MA10), with every $2m$ replaced by $2m+\Lambda$,
gives
\begin{equation}\tag{MA11}\begin{aligned}
 \int_\R(1-e^{-\Lambda|v|})FK\dd v
 ={}&C\sum_{m\ge0}\left[
 \frac1{2m+w}-\frac1{2m+w+\Lambda}
 +\frac1{2m+z}-\frac1{2m+z+\Lambda}\right]\\
 &-2a x^{-1/2}\sum_{m\ge0}
 \frac{x^{-2m}}{(2m+z)(2m+p)}\\
 &+2a x^{-1/2-\Lambda}\sum_{m\ge0}
 \frac{x^{-2m}}{(2m+z+\Lambda)(2m+p+\Lambda)}.
\end{aligned}\end{equation}
Each bracketed harmonic difference is $O_{\Lambda,\sigma_0,t}(m^{-2})$;
the two rational series are separately absolutely convergent. Thus (MA11)
splits only convergent series, not divergent harmonic sums.

\section{The exact digamma difference and its complex limit}
Euler's locally uniformly convergent product for $1/\Gamma$, as recorded in
DLMF 5.8.2 \cite{Gamma}, gives by local logarithmic differentiation
\begin{equation}\tag{MA12}
 \psi(b)-\psi(a_0)
 =\sum_{m=0}^\infty\left(\frac1{m+a_0}-\frac1{m+b}\right),
\end{equation}
for $a_0,b$ away from the nonpositive integers. The differentiated differences
converge locally uniformly since they are $O(m^{-2})$. One can first
differentiate on the positive real half-plane and then continue meromorphically;
this yields the same expression without imposing a logarithm branch at
$z/2$, whose real part lies in $(-1/2,0)$.
Applying (MA12) to (MA11) evaluates its first line exactly as
\begin{equation}\tag{MA13}
 \frac C2\left[
 \psi((w+\Lambda)/2)-\psi(w/2)
 +\psi((z+\Lambda)/2)-\psi(z/2)\right].
\end{equation}

For clarity, the limiting formula needed here follows directly from the same
product. For real $r>0$ it gives
$\psi(r)=\lim_{N\to\infty}[\log N-\sum_{m=0}^N(m+r)^{-1}]$.
The decreasing-integrand sum comparison gives
\[
 0\le\sum_{m=0}^N\frac1{m+r}
      -\int_0^{N+1}\frac{du}{u+r}\le\frac1r,
 \qquad |\psi(r)-\log r|\le\frac1r.
\]
For $r\ge\max(1,2|c_0|)$, (MA12) also gives
\[
 |\psi(r+c_0)-\psi(r)|
 \le2|c_0|\sum_{m\ge0}(m+r)^{-2}
 \le\frac{4|c_0|}{r}.
\]
Taking $r=\Lambda/2$, $c_0=s/2$ proves the complex estimate
\begin{equation}\tag{MA14}
 \left|\psi((\Lambda+s)/2)-\log\Lambda+\log2\right|
 \le\frac{2+4|s|}{\Lambda},\qquad
 \Lambda\ge\max(2,2|s|).
\end{equation}
Thus no passage to real parts of the complex identity is involved.
The shifted rational series in (MA11) satisfies, for $\Lambda>2$,
\begin{equation}\tag{MA15}
 \left|\sum_{m\ge0}
 \frac{x^{-2m}}{(2m+z+\Lambda)(2m+p+\Lambda)}\right|
 \le\frac1{(\Lambda-1)^2}+\frac1{2(\Lambda-1)}.
\end{equation}
Indeed the real parts of the two denominator factors exceed
$2m+\Lambda-1$, and the last bound is the sum-plus-integral bound for
$(2u+\Lambda-1)^{-2}$. Multiplication by $x^{-\Lambda}$ makes the error
smaller when $x>1$; (MA15) also proves its convergence to zero directly
at $x=1$.

Subtracting $C\log\Lambda$ in (MA11), then using (MA13)--(MA15), proves
\begin{equation}\tag{MA16}\boxed{\begin{aligned}
 \PF(FK)={}&-C\log2
 -\frac C2[\psi(w/2)+\psi(z/2)]\\
 &-2a x^{-1/2}\sum_{m=0}^\infty
       \frac{x^{-2m}}{(2m+z)(2m+p)}.
\end{aligned}}\end{equation}
More quantitatively, for
$\Lambda\ge\max(3,2|w|,2|z|)$ the difference between the regulated finite
part in (MA2) and (MA16) is bounded by
\begin{equation}\tag{MA17}
 \frac{|C|[2+2(|w|+|z|)]}{\Lambda}
 +2a x^{-1/2-\Lambda}
 \left[\frac1{(\Lambda-1)^2}+\frac1{2(\Lambda-1)}\right].
\end{equation}
This controls the actual regulator with its original subtraction.
For $M\ge1$, the tail of the unshifted rational series after $m=M$ satisfies
\begin{equation}\tag{MA18}
 \left|\sum_{m>M}\frac{x^{-2m}}{(2m+z)(2m+p)}\right|
 \le\frac1{2M}.
\end{equation}
For $x>1$ it also satisfies the stronger bound
$x^{-2M-2}/[2(M+1)^2(1-x^{-2})]$.
To verify this, use $|2m+z|>2m-1\ge m$ and $|2m+p|>2m$ for $m\ge1$,
then the integral or geometric bound. The separate $m=0$ term is finite
since $z,p\ne0$.

\section{Both prime orientations and the exact reflection identity}
Let $\Lambda(n)$ be the von Mangoldt function, with $\Lambda(1)=0$.
The two orientations are, directly from (MA9),
\begin{equation}\tag{MA19}\begin{aligned}
 n^{-1/2}F(\log n)
 &=\begin{cases}
 Cn^{\sigma_0-1-it},& n\le x,\\
 x^{\sigma_0-1/2+it}n^{-\sigma_0-it},& n>x,
 \end{cases}\\
 n^{-1/2}F(-\log n)&=Cn^{-w}.
\end{aligned}\end{equation}
At $n=x$ the two displayed positive-orientation expressions both equal
$n^{-1/2}$. Thus the split does not create a jump or a half-weight at
a prime power. At $x=1$ the first sum contains only $n=1$ and is zero.
Absolute convergence of the infinite prime sums follows from
$0\le\Lambda(n)\le\log n$ and $\sigma_0>1$.

For completeness, the Euler product for $\zeta(s)$ on $\Re s>1$
\cite{Zeta} is nonzero and has the locally absolutely convergent logarithm
$\sum_{\ell\ \mathrm{prime}}\sum_{r\ge1}\ell^{-rs}/r$.
Absolute convergence follows from
$\sum_\ell\sum_{r\ge1}\ell^{-r\sigma_0}/r<\infty$;
local termwise differentiation is permitted by the corresponding convergent
$\log n\,n^{-\sigma_0}$ bound. Unique prime-power factorization therefore
gives the exact signs
\begin{equation}\tag{MA20}
 P_-:=\sum_{n\ge1}\Lambda(n)n^{-1/2}F(-\log n)
 =C\sum_{n\ge1}\Lambda(n)n^{-w}
 =-C\frac{\zeta'(w)}{\zeta(w)},\qquad
 -P_-=+C\frac{\zeta'(w)}{\zeta(w)}.
\end{equation}
For example, for any integer $N$ beyond $e^{1/\sigma_0}$, the absolute tail
of the Dirichlet sum after $N$ is at most
\[
 \int_N^\infty(\log u)u^{-\sigma_0}\dd u
 =N^{1-\sigma_0}\left[\frac{\log N}{\sigma_0-1}
                       +\frac1{(\sigma_0-1)^2}\right].
\]
Its outside factor $|C|$ is retained for $P_-$.

Put $\mathscr L(s)=\pi^{-s/2}\Gamma(s/2)\zeta(s)$.
The exact zeta functional equation is $\mathscr L(s)=\mathscr L(1-s)$,
equivalently the $\xi$ reflection with its $s(s-1)/2$ factor
\cite{Zeta}. At $w$ the Euler product and the nonvanishing Gamma factor
give $\mathscr L(w)\ne0$. At $z=1-w$, which has $-1<\Re z<0$,
$\Gamma(z/2)$ is finite and nonzero, so reflection also proves
$\zeta(z)\ne0$. Thus every logarithmic derivative below exists.
Differentiating the equality as a meromorphic equality and dividing by its
nonzero value gives
\begin{equation}\tag{MA21}\begin{aligned}
 -\tfrac12\log\pi+\tfrac12\psi(z/2)+\frac{\zeta'(z)}{\zeta(z)}
 &=-\left[-\tfrac12\log\pi+\tfrac12\psi(w/2)
                           +\frac{\zeta'(w)}{\zeta(w)}\right],\\
 \frac{\zeta'(w)}{\zeta(w)}-\log\pi
 +\tfrac12[\psi(w/2)+\psi(z/2)]
 &=-\frac{\zeta'(z)}{\zeta(z)}.
\end{aligned}\end{equation}
This derivation uses ratios, not an unproved global logarithm branch.

\section{Pole cancellation and the completed identity}
Substitute (MA16) and (MA20) into the literal residual expression. Its
constant is $\log A=-\log(2\pi)$, so $-\log(2\pi)+\log2=-\log\pi$.
Using (MA21) only after that cancellation yields
\begin{equation}\tag{MA22}\begin{aligned}
 &\Phi(0)+C\log A-\PF(FK)-P_-\\
 &=\Phi(0)+C\left\{\frac{\zeta'(w)}{\zeta(w)}-\log\pi
              +\frac{\psi(w/2)+\psi(z/2)}2\right\}
       +2a x^{-1/2}\sum_{m\ge0}\frac{x^{-2m}}{(2m+z)(2m+p)}\\
 &=-C\frac{\zeta'(z)}{\zeta(z)}
   +\underbrace{\Phi(0)+\frac{2a x^{-1/2}}{zp}}_{=0\text{ by (MA4)}}
   +2a x^{-1/2}\sum_{m\ge1}\frac{x^{-2m}}{(2m+z)(2m+p)}.
\end{aligned}\end{equation}
The two denominator presentations are exactly related by
\[
 2m+z=-(\sigma_0-1-it-2m),\qquad 2m+p=\sigma_0+it+2m.
\]
Therefore (MA22) is precisely the required complete scalar identity
\begin{theorem}
For every original parameter tuple (MA1), with the literal finite part
(MA2), all the terms below exist and
\begin{equation}\tag{MA23}\begin{aligned}
 &\Phi(0)+C\log((2\pi)^{-1})-\PF(FK)
   -\sum_{n\ge1}\Lambda(n)n^{-1/2}F(-\log n)\\
 &\quad=-C\frac{\zeta'(1-\sigma_0+it)}{\zeta(1-\sigma_0+it)}
 -x^{-1/2}\sum_{m\ge1}
 \frac{(2\sigma_0-1)x^{-2m}}
      {(\sigma_0-1-it-2m)(\sigma_0+it+2m)}.
\end{aligned}\end{equation}
\end{theorem}
The proof has not used a zero sum, a zero-location hypothesis, or the
incorrect printed Guinand--Weil residual equality.

\section{The endpoint and the receiving Montgomery formula}
At $x=1$, the actual test is $F(v)=e^{-a|v|}e^{-itv}$ and $C=1$.
The direct two-interval integration was already valid with $U=0$.
Independently, (MA12) and $p-z=2a$ give
\[
 2a\sum_{m\ge0}\frac1{(2m+z)(2m+p)}
 =\sum_{m\ge0}\left(\frac1{2m+z}-\frac1{2m+p}\right)
 =\tfrac12[\psi(p/2)-\psi(z/2)].
\]
Substitution in (MA16) proves the endpoint value
\begin{equation}\tag{MA24}
 \boxed{\PF(FK)=-\log2-\tfrac12
 [\psi((\sigma_0-it)/2)+\psi((\sigma_0+it)/2)]\quad(x=1).}
\end{equation}
All series converge directly at this endpoint by (MA18); no omitted
limiting or prime-power convention is required.

We now receive the established identity into the recorded original Weil
formula, with its symmetric zero cutoff. Its rational trivial-character
specialization is \cite[formula (11)]{Weil}
\begin{equation}\tag{MA25}
 \lim_{T\to\infty}\sum_{|\Im\rho|<T}\Phi(\rho)
 =\Phi(0)+\Phi(1)+C\log A-P_+-P_--\PF(FK),
\end{equation}
where $P_+$ is the positive-orientation sum from (MA19) and all nontrivial
zeros are counted with multiplicity. We use this classical explicit formula
at its original test domain, whose hypotheses were verified after (MA2).
The full source provides its literal normalization and the two prime orientations.
Substitution of the independently proved (MA23) gives the zero-side formula
without changing any zero coordinate.

On the real-ordinate domain used in Montgomery's 1973 paper, namely RH,
$\rho=1/2+i\gamma$ with $\gamma\in\R$ and
\[
 \Phi(\rho)=\frac{(2\sigma_0-1)x^{i\gamma}}
 { (\sigma_0-1/2)^2+(t-\gamma)^2}.
\]
Equations (MA4), (MA19), (MA23) and (MA25) therefore produce exactly
\begin{equation}\tag{MA26}\begin{aligned}
 &(2\sigma_0-1)\sum_\gamma
 \frac{x^{i\gamma}}{(\sigma_0-1/2)^2+(t-\gamma)^2}\\
 &=-x^{-1/2}\left[
 \sum_{n\le x}\Lambda(n)(x/n)^{1-\sigma_0+it}
 +\sum_{n>x}\Lambda(n)(x/n)^{\sigma_0+it}\right]\\
 &\quad-C\frac{\zeta'(1-\sigma_0+it)}{\zeta(1-\sigma_0+it)}
 +\frac{x^{1/2}(2\sigma_0-1)}{(\sigma_0-1+it)(\sigma_0-it)}\\
 &\quad-x^{-1/2}\sum_{m\ge1}
 \frac{(2\sigma_0-1)x^{-2m}}
 {(\sigma_0-1-it-2m)(\sigma_0+it+2m)}.
\end{aligned}\end{equation}
This is Montgomery's complete retained equation \cite[pp.~185--187]{Montgomery}.
The sum means at least the original symmetric cutoff in (MA25); its
convergence is supplied by that cited theorem. No later real-ordinate estimate
is declared valid for a complex ordinate by this receiving calculation.
The scalar (MA23) itself is unconditional on (MA1); the RH qualification in
(MA26) belongs only to the original choice of real zero coordinates.

\section{Exact backward replacement}
In the complete Original Six Readers provider, the proved transform,
prime-range and pole formulas remain as printed. Replace the status sentence
``The remaining gamma/trivial-zero commutator is not claimed proved'' and
the concluding paragraph beginning ``Still to be proved is'' in
\texttt{proofs/B\_guinand\_weil\_montgomery\_interfaces.tex} by (MA23),
with the proof (MA5)--(MA24). In
\texttt{satellites/10\_literature\_foundations.tex}, replace the paragraph
beginning ``What is still open is one explicit commutator'' by the same
completed identity and its exact parameter domain. These are successor
replacement instructions; no archived source bytes are changed here.

The separate local labels \texttt{eq:guinand-weil-arch-map} and
\texttt{eq:guinand-weil-residual-complete} need the nonzero correction proved
in the companion note, not (MA23) with a discarded term. On its explicit
common test subspace the exact additional functional is
\[
 \mathcal D(f)=\int_0^\infty f(v)
       [2e^{-v/2}-\operatorname{sech}(v/2)]\dd v.
\]
The original right side must receive $+\mathcal D(f)$. In particular this
change must be propagated to the asserted local subtraction proof and its
status records; it is not a correction to Guinand's published theorem.
The companion proves the exact test-space inclusion with decay exponent
strictly greater than $1/2$, all jump terms, and a nonzero admissible example.

\begin{thebibliography}{9}
\bibitem{Gamma} R. A. Askey and R. Roy, \emph{Gamma Function}, NIST Digital
Library of Mathematical Functions, Chapter 5, equation 5.8.2 (Euler product)
and \S5.9(ii) (digamma differences).
\url{https://dlmf.nist.gov/5.8.E2}; \url{https://dlmf.nist.gov/5.9}.
The required difference and its complex large-argument limit are proved
directly in (MA12)--(MA14).
\bibitem{Zeta} T. M. Apostol, \emph{Zeta and Related Functions}, NIST DLMF,
Chapter 25, \S25.2(iv) (Euler product), equations 25.4.3--4 (completed
reflection and the exact $\xi$ convention).
\url{https://dlmf.nist.gov/25.2}; \url{https://dlmf.nist.gov/25.4}.
\bibitem{Weil} A. Weil, \emph{Sur les formules explicites de la th\'eorie
des nombres premiers}, Communications du S\'eminaire Math\'ematique de
l'Universit\'e de Lund, supplementary volume dedicated to Marcel Riesz
(1952), 252--265; reprinted in \emph{Collected Papers}, vol.~II,
Springer, 1979, 48--61, conditions (A),(B), finite part and formula (11),
52--58. The complete original Six Readers provider records these exact
conventions rather than substituting another finite part.
\bibitem{Montgomery} H. L. Montgomery, \emph{The pair correlation of zeros
of the zeta function}, Analytic Number Theory (St.~Louis, 1972),
Proceedings of Symposia in Pure Mathematics 24, American Mathematical
Society (1973), 181--193, equations (18)--(22).
\url{https://doi.org/10.1090/PSPUM/024/9944}; author's primary paper:
\url{https://websites.umich.edu/~hlm/paircor1.pdf}.
\end{thebibliography}
\end{document}
